

%% ===================================================================
\section{선행연구 검토 및 연구의 차별점}
\label{sec:lit-gap}
%% ===================================================================

앞에서 검토한 BCM, 정적 RTO의 한계, 관측성, 불확실성, 시그모이드 바운딩의 이론적 논의를 종합하고 선행연구의 흐름을 정리하여 본 연구가 수행하고자 하는 차별점을 [표 2-8]에 주제별로 정리하였다.

%% -------------------------------------------------------------------
%% -------------------------------------------------------------------

\begin{table}[htbp]
\centering
\caption{주제별 선행연구 요약}
\label{tab:prior-research}
\small
\begin{tabularx}{\textwidth}{@{}>{\raggedright\arraybackslash}p{1.8cm}
  >{\raggedright\arraybackslash}p{3.0cm}
  >{\raggedright\arraybackslash}X
  >{\raggedright\arraybackslash}p{4.0cm}@{}}
\toprule
\textbf{주제} & \textbf{주요 문헌} & \textbf{핵심 내용} & \textbf{한계} \\
\midrule
BCM/BCMS 표준 &
\citet{herbane2010disaster}, ISO 22301(2019) &
PDCA 기반 BCMS 프레임워크, BCM 발전 3단계 분류 &
RTO를 정적 고정값으로 설정, 동적 조정 메커니즘 부재 \\
\addlinespace
BIA 및 RTO &
\citet{snedaker2013bcdr}, \citet{hiles2010bcm} &
BIA 방법론 체계화, MTPD-RTO 관계 정립 &
RTO 동적 특성 미논의, 관측성 미반영 \\
\addlinespace
\m{관측성(제어 이론)} &
\citet{kalman1960control}, \citet{avizienis2004basic} &
LTI 시스템 관측성 형식화, 의존성 분류 체계 &
BCM 영역 적용 부재, RTO와의 수리적 연결 미시도 \\
\addlinespace
\m{관측성(IT 운영)} &
\citet{beyer2016sre}, \citet{majors2022observability} &
SRE 관측성 실무, 세 가지 핵심 요소(three pillars) 체계 &
정량적 모델화 미수행, MTTD/MTTR 관계만 정성적 \\
\addlinespace
가우시안 불확실성 &
\citet{jaynes2003probability}, \citet{gelman2013bayesian} &
CLT 기반 불확실성 정량화, 베이지안 추론 체계 &
극단값 과소평가, 파레토 현상 미포착 \\
\addlinespace
파레토 분포/EVT &
\citet{taleb2007blackswan}, \citet{embrechts1997extremal}, \citet{newman2005powerlaws} &
블랙 스완 개념, EVT 수학적 기초, 멱법칙 보편성 &
금융/보험 중심 적용, BCM/RTO 영역 미적용 \\
\addlinespace
국내 BCM &
\citet{cheung2019quickhit}, \citet{jang2021bcms}, \citet{jung2023bcms} &
BIA Quick-Hit 프레임워크, BCMS 성과 실증, BCM 평가지표 개발 &
정적 RTO 전제, 동적 조정 모델 미제시 \\
\addlinespace
DRTO, 시그모이드 바운딩 &
\citet{lee2026drto} &
관측성 기반 DRTO에 시그모이드 바운딩 도입, 곡률 매개변수 $\kappa^*$ 도출 &
관측성 단일 채널 한정, 불확실성 채널 미포함 \\
\bottomrule
\end{tabularx}
\end{table}

[표 2-8]에서 확인되는 선행연구의 공통적 한계는 \chg{다음 세 가지로 요약되며 2.7(DRTO 모델 구조)에서 상세히 다룬다. 첫째는 관측성을 통합한 RTO 모델의 부재로서 부분적인 단일 채널
모델만 존재하고~\citep{lee2026drto}, 둘째는 불확실성이 정량화된 동적 RTO 모델의 부재이며,
셋째는 극단 시나리오에 대한 파레토 분석의 부재이다.} 본 연구의 $\DRTO$ 프레임워크는 이
\chg{세 가지 한계}를 통합적으로 해소하여 선행연구와 차별화하는 것을 목표로 한다.


%% -------------------------------------------------------------------
%% -------------------------------------------------------------------

\p{[}표~\ref{tab:prior-research}\p{]}의 선행연구를 종합하면, 기존 BCM/RTO 연구와
본 연구의 $\DRTO$ 프레임워크 사이에는 다음 세 가지 핵심적 차별점이 존재한다.

\chg{첫째, 관측성을 통합한 RTO 모델의 부재이다.} 기존 BCM 연구에서 RTO는 BIA를 통해 고정적으로 결정되며, 시스템의 관측성 수준이 복구 시간에 미치는 영향은 정량적으로 모델링되지 않았다. 제어 이론\citep{kalman1960control}과 SRE\citep{beyer2016sre, majors2022observability}에서 관측성의 중요성이 \m{인정되며}, \citet{lee2026drto}은 관측성 단일 채널을 시그모이드 바운딩으로 RTO에 연결하는 수리적 프레임워크를 처음 제시하였으나, 불확실성 채널과 결합한 다채널 통합 모델은 여전히 부재하였다.

\chg{둘째, 불확실성이 정량화된 동적 RTO($\DRTO$) 모델의 부재이다.} 기존 RTO는 결정론적(deterministic) 단일값으로 설정되어 복구 과정에서 발생하는 불확실성을 반영하지 못하였다. 확률론적 사고방식\citep{gelman2013bayesian, jaynes2003probability}이 위험관리 분야에 일부 \m{적용되나}, RTO 자체를 확률 분포로 모델링하여 불확실성에 따른 동적 조정을 수행하는 연구는 부재하였다.

\chg{셋째, 극단적 시나리오에 대한 파레토 분석의 부재이다.} \m{극단값이론(EVT)}과 파레토 분포는 금융 위험관리\citep{embrechts1997extremal, mcneil2015risk}에서 활발히 \m{활용되나}, BCM 영역에서 극단적 복구 지연 시나리오를 파레토 분포로 모델링한 연구는 찾아보기 어렵다. \citet{taleb2007blackswan}이 지적한 극단 사건의 과소평가 문제가 RTO 설정에서도 동일하게 \m{발생함에도} 불구하고, 이에 대한 체계적 분석 및 대응 방안은 선행 연구에서 다루어지지 않았다.

[표 2-9]는 이상의 논의를 종합하여 기존 연구의 접근 방식과 본 연구의 DRTO 프레임워크가 제시하는 해결 방안을 대비하여 비교한 것이다.

\begin{table}[htbp]
\centering
\caption{선행 연구의 한계와 본 연구의 접근 비교}
\label{tab:research-gap}
\small
\begin{tabularx}{\textwidth}{@{}>{\raggedright\arraybackslash}p{2.6cm}
  >{\raggedright\arraybackslash}X
  >{\raggedright\arraybackslash}X
  >{\raggedright\arraybackslash}p{2.3cm}@{}}
\toprule
\textbf{연구 차별점} & \textbf{기존 연구} & \textbf{본 \m{연구($\DRTO$)}} & \textbf{관련 조건} \\
\midrule
차별점 1:\newline 관측성 미반영 &
단일 채널 관측성 모델~\citep{lee2026drto}만 존재;\newline 다채널 관측성-불확실성\newline 통합 부재 &
$O_{\text{raw}} \in [0,1]$ 관측성 원시값을\newline 개별 시그모이드 바운딩의\newline 핵심 입력으로 통합 &
\p{C1 (관측성)} \\
\addlinespace
차별점 2:\newline 불확실성 미정량 &
결정론적 단일값 RTO;\newline 확률적 변동 미고려 &
가우시안 $N(\mu,\sigma^2)$ 및\newline 파레토 $\text{Pareto}(\alpha)$\newline 이중채널 모델링 &
\p{C2 (가우시안),\newline C3 (파레토)} \\
\addlinespace
차별점 3:\newline 극단 시나리오\newline 미분석 &
BCM 영역에서 EVT/\newline 파레토 미적용;\newline 꼬리 위험 미포착 &
EVT 기반 파레토 분석;\newline 분할 회귀(Core/Tail);\newline 이중채널 비선형 효과 정량화 &
\p{C3 tail (파레토\newline P85.8 이상)} \\
\addlinespace
바운딩 기법 &
클리핑(hard boundary)\newline 또는 미적용 &
개별 시그모이드 바운딩\newline $S(z) = 2/(1+e^{-z})-1$\newline $\kappa = 1.2$, 전 구간 미분 가능 &
전 조건 공통 \\
\bottomrule
\end{tabularx}
\end{table}

\chg{[표 2-9]의 관련 조건 열에 나타나는 조건 표기는 관측성과 불확실성의 활성화 조합에 따라 정의한 네 가지 실험 조건을 가리킨다. C0는 관측성과 불확실성을 모두 비활성화한 기준선이고, C1은 관측성만 활성화한 조건이며, C2는 C1에 가우시안 불확실성을 더한 조건, C3는 C1에 파레토 불확실성을 더한 조건이다.}

본 연구의 $\DRTO$ 프레임워크는 위의 세 가지 차별점을 통합적으로 구현하는 것을 목표로 한다.
구체적으로, 관측성을 핵심 입력 변수로 도입하고,
가우시안 및 파레토 분포를 통해 불확실성을 이중 채널로 정량화하며,
파레토 분포를 BCM 영역에 도입하여 극단 시나리오 분석의 공백을 채운다.
다음으로 이러한 이론적 기반 위에 구축되는 DRTO 프레임워크의 구체적 수리 모델과 가설 체계를 상세히 기술한다.
